%% 第 23 章：正算子与正平方根
%% 本文件由 tools/md2tex.py 自动生成，请勿直接编辑。

\chapter{正算子与正平方根}


\begin{jdefn}{正算子的定义}{r:24:1}
若 $T$ 自伴，且：

\[
\langle Tv,v\rangle\ge0
\]

对所有 $v\in V$ 成立，则称 $T$ 是正算子。

正算子描述一种不会在任何方向上产生负二次能量的变换。这便于我们所谓的二次型产生联系了。
\end{jdefn}

\begin{jprop}{正算子的谱刻画}{r:24:2}
\[
T\text{ 正}
\iff
T\text{ 自伴，且所有特征值非负}.
\]
\end{jprop}

\begin{jproof}{证明思路}
若：

\[
Te_j=\lambda_je_j,
\]

则：

\[
\langle Te_j,e_j\rangle
=
\lambda_j.
\]

所以正性要求每个特征值都非负。

反过来，在规范正交特征基下：

\[
v=\sum_j a_je_j.
\]

于是：

\[
\langle Tv,v\rangle
=
\sum_j\lambda_j|a_j|^2.
\]

若全部 $\lambda_j\ge0$，该和必然非负。

又因为前面证明过，在 $\mathbb C^n$ 上只有自伴算子才能使二次型都是实数，所以，正算子是自伴的。
\end{jproof}

\begin{jthm}{正平方根}{r:24:3}
若 $T$ 是正算子，则存在唯一的正算子 $R$，使：

\[
R^2=T.
\]

记作：

\[
R=T^{1/2}.
\]
\end{jthm}

\section{正平方根的构造}

在 $T$ 的规范正交特征基下：

\[
Te_j=\lambda_je_j,
\qquad
\lambda_j\ge0.
\]

定义：

\[
Re_j=\sqrt{\lambda_j}e_j.
\]

则：

\[
R^2e_j=\lambda_je_j=Te_j.
\]

因此：

\[
R^2=T.
\]

\section{正平方根的唯一性}

正平方根在每个特征方向上的特征值必须是非负数：

\[
\sqrt{\lambda_j}.
\]

所以它被 $T$ 的特征结构唯一确定。（由前面的定理，线性映射在给定在一组基上的取值时候，这个映射便被唯一确定了)

注意：普通平方根未必唯一。例如恒等算子在不同特征方向上可以取 $1$ 或 $-1$；但其中只有全取正的那个平方根是正平方根。
